\section{本章小结}
\label{sec:wavelet-summary}

本章要找一种坐标，让平稳区域的细节接近零，让局部跳变只影响少数系数。小波系数是把函数这条无限维向量换到多尺度标准正交基后读出的坐标，它按尺度$j$和位置$k$排成一张数表。

有限Haar变换给出最简单的回答：对$2^J$个取值反复求平均与半差，得到一个总平均和$2^J-1$个按尺度、位置排列的差值，一处改动在每一层只影响一个系数。把平均与半差乘以$\sqrt2$，变换成为正交矩阵，矩阵的各行就是新的基向量，系数的平方和等于原向量的平方和，逆变换就是转置。

放回区间上，有限变换就是在分段常数函数空间$V_J$中更换标准正交基。到$V_j$的投影是取区间平均；细节空间$W_j=V_{j+1}\cap V_j^{\perp}$由Haar函数张成，记录较粗表示中缺少的左右差异。取值向量要乘以$2^{-J/2}$才是$L^2$坐标，此时保持的是函数的平方积分。

二进阶梯函数在$L^2([0,1])$中稠密，所以粗项加上全部细节构成标准正交基，部分和就是区间平均$P_Jf$，按$L^2$收敛并满足Parseval等式。区间上的粗项不能省略；连续函数的部分和一致收敛，跳跃函数的部分和不一致收敛。

实轴上，正交多分辨率分析把这些结构抽象成嵌套、缩放、平移基、稠密与分离几条要求。尺度函数通过二尺度关系连接相邻层，母小波由二尺度系数写出，证明的关键是把序列按奇偶分成两路，在每个频率上得到$\mathbb C^2$的一组标准正交基。实轴上每层有无限多个位置，粗项在起始层趋于$-\infty$时消失。

消失矩让局部的低次多项式不留痕迹，光滑函数的细节系数因而随层数快速衰减：对$x^2$，Haar的系数是$2^{-3j/2}$的阶，D4的系数精确等于$-(\sqrt3/8)2^{-5j/2}$。小波的光滑性则决定有限项重构是否光滑。

有限区间要处理跨越端点的基函数。周期化给出周期函数空间的标准正交基，零阶矩保留而高阶矩在跨接缝处失效。周期延拓是$L^2$等距同构，却可能在接缝处产生跳跃或折角；反射延拓能消除取值跳跃而不保证导数连续；边界小波在区间内重新构造端点附近的基函数，需要重新核对正交、完备与消失矩。

与Fourier系数比较，全局振荡适合频率坐标，局部跳跃适合尺度位置坐标，哪种误差更小取决于函数、基和预算的计法。截断误差的定量刻画在第~\ref{chap:function-approximation}章研究；第~\ref{chap:besov-spaces}章在本章的周期边界设定下，用小波系数在各层上的大小定义和刻画函数的正则性。
